\renewcommand{\Anno}{2026}

\renewcommand{\NumeroProblema}{1}

\renewcommand{\Testo}{
Tre ingegneri, Marco, Pasquale e Thomas, devono calcolare l'area di un quadrilatero, di cui conoscono le lunghezze dei lati. Non sapendo come fare, provano a trovare la risposta con metodi creativi.
\begin{itemize}
    \item Marco usa la media dei quadrati delle lunghezze dei lati.
    \item Pasquale moltiplica tra di loro le medie delle due coppie di lati opposti.
    \item Thomas sceglie un vertice del quadrilatero, e poi moltiplica tra di loro la media dei due lati che hanno quel vertice come estremo e la media dei restanti due lati.
\end{itemize}

% SSSUP-FIGURA-RICOSTRUITA: es1.pdf, schema redazionale PSTricks.
% Fonte di confronto: fonti-locali/2026/es1.pdf (conservata senza modifiche).
% Le coordinate sono illustrative; i numeri sui lati sono i dati del problema.
\begin{center}
\newrgbcolor{SSSUPGriglia}{0.30 0.78 0.77}
\psset{unit=0.88cm}
\begin{pspicture}(-0.62,-0.45)(4.54,5.60)
  \multido{\nx=0+1}{5}{%
    \psline[linecolor=SSSUPGriglia,linewidth=0.4pt](\nx,0)(\nx,5)%
  }
  \multido{\ny=0+1}{6}{%
    \psline[linecolor=SSSUPGriglia,linewidth=0.4pt](0,\ny)(4,\ny)%
  }
  \pspolygon[fillstyle=solid,fillcolor=lightgray,linewidth=1.4pt]%
    (0,2)(0,4)(4,5)(2.5,0.25)
  \rput[r](-0.10,2){$A$}
  \rput[r](-0.10,4){$B$}
  \rput[bl](4.04,5.00){$C$}
  \rput[t](2.58,0.17){$D$}
  \rput[r](-0.18,3){$2$}
  \rput[b](1.8,5.12){$4$}
  \rput[l](3.38,2.66){$5$}
  \rput[b](1.43,0.82){$3$}
\end{pspicture}
\end{center}

Per esempio, se le lunghezze dei lati, presi in senso orario a partire dal vertice scelto da Thomas, sono 2, 4, 5, 3 (come mostrato in figura a partire dal vertice A), Marco, Pasquale e Thomas rispondono, rispettivamente,
$$
M=\frac{2^2+4^2+5^2+3^2}{4},
\qquad\qquad
P=\frac{2+5}{2}\cdot\frac{4+3}{2},
\qquad\qquad
T=\frac{2+3}{2}\cdot\frac{4+5}{2}.
$$

\begin{enumerate}
\renewcommand{\labelenumi}{(\alph{enumi})}

\item Determinare per quali quadrilateri Marco ottiene la risposta corretta.

\item Determinare per quali quadrilateri Pasquale ottiene la risposta corretta.

\item Determinare per quali quadrilateri il Thomas può ottenere la risposta corretta scegliendo opportunamente il vertice.

\end{enumerate}
}

\renewcommand{\Risposta}{
% SSSUP-BOZZA-AI: Risposta; autore=ChatGPT; data=2026-11-08; stato=DA-VERIFICARE
% SSSUP-PROVENIENZA: AI: rielaborazione e sviluppo di appunti forniti
Per quadrilateri semplici non degeneri:
\begin{enumerate}
\item[(a)] Marco ottiene l'area corretta esattamente per i quadrati.
\item[(b)] Pasquale ottiene l'area corretta esattamente per i rettangoli.
\item[(c)] Thomas può ottenerla esattamente per gli aquiloni ciclici: unione di due triangoli rettangoli congruenti lungo l'ipotenusa comune. Deve scegliere uno degli estremi di tale ipotenusa, dove si incontrano due lati uguali.
\end{enumerate}
}

\renewcommand{\Aiutino}{
% SSSUP-BOZZA-AI: Aiutino; autore=ChatGPT; data=2026-11-08; stato=DA-VERIFICARE
% SSSUP-PROVENIENZA: AI: rielaborazione e sviluppo di appunti forniti
L'area di un triangolo non supera la metà del prodotto di due suoi lati. Applica questa stima tagliando il quadrilatero secondo entrambe le diagonali. Per Thomas, usa inoltre $2K\leq AC\cdot BD\leq AB\cdot CD+BC\cdot DA$ e analizza tutti i casi di uguaglianza.
}

\renewcommand{\SoluzioneUfficiale}{}

\renewcommand{\SoluzioniAlternative}{}

\renewcommand{\Osservazioni}{}

\renewcommand{\Commenti}{}

% Aggiunta redazionale: la soluzione preesistente non viene modificata.
\providecommand{\SoluzioneRedazionale}{}
\renewcommand{\SoluzioneRedazionale}{
% SSSUP-BOZZA-AI: SoluzioneRedazionale; autore=ChatGPT; data=2026-11-08; stato=DA-VERIFICARE
% SSSUP-PROVENIENZA: AI: rielaborazione e sviluppo di appunti forniti
La soluzione rielabora la bozza fornita. Consideriamo un quadrilatero semplice non degenere $ABCD$, con lati consecutivi $a,b,c,d$ e area $K$. Anche nel caso concavo, l'area è non superiore alla somma delle aree dei due triangoli determinati da una diagonale. Quindi
\[
2K\leq ab+cd,\qquad2K\leq ad+bc.
\]
L'uguaglianza nella prima stima richiede una decomposizione senza sovrapposizioni e angoli retti in $B,D$; nella seconda richiede angoli retti in $A,C$. I quadrilateri di uguaglianza risultano dunque convessi.

\paragraph{Punto (a).}
Usando $2ab\leq a^2+b^2$ e $2cd\leq c^2+d^2$,
\[
K\leq\frac{ab+cd}{2}\leq\frac{a^2+b^2+c^2+d^2}{4}=M.
\]
L'uguaglianza richiede $B,D$ retti, $a=b$ e $c=d$. I due triangoli rettangoli isosceli hanno la medesima ipotenusa $AC$, quindi sono congruenti e tutti e quattro i lati sono uguali. Il quadrilatero è un quadrato; viceversa ogni quadrato realizza l'uguaglianza.

\paragraph{Punto (b).}
Sommando le due stime iniziali,
\[
4K\leq ab+cd+ad+bc=(a+c)(b+d).
\]
L'uguaglianza richiede tutti e quattro gli angoli retti, cioè un rettangolo. Per un rettangolo la formula di Pasquale coincide effettivamente con l'area.

\paragraph{Punto (c).}
Supponiamo che Thomas scelga $A$. Per le diagonali vale
\[
2K\leq AC\cdot BD\leq ac+bd.
\]
La prima disuguaglianza segue dalla formula dell'area tramite il prodotto vettoriale delle diagonali; la seconda è la disuguaglianza di Tolomeo. Per completezza, identificando il piano con i complessi, quest'ultima segue dalla disuguaglianza triangolare applicata all'identità
\[
(z_A-z_C)(z_B-z_D)=(z_A-z_B)(z_C-z_D)+(z_A-z_D)(z_B-z_C).
\]
Sommando con $2K\leq ab+cd$ otteniamo
\[
4K\leq ab+cd+ac+bd=(a+d)(b+c),
\]
ossia $K\leq T$.

Per avere uguaglianza sono necessari angoli retti in $B,D$ e diagonali perpendicolari. I due angoli retti collocano $B,D$ sulla circonferenza di diametro $AC$. Poiché il diametro $AC$ è perpendicolare alla corda $BD$, ne è l'asse: $B$ e $D$ sono simmetrici rispetto ad $AC$. Quindi $AB=AD$ e $BC=CD$, e i due triangoli rettangoli sono congruenti.

Viceversa, incollando due triangoli rettangoli congruenti lungo l'ipotenusa si ottiene un aquilone ciclico. Scegliendo come $A$ uno degli estremi dell'ipotenusa, si ha $a=d$, $b=c$, $K=ab$ e $T=ab$. Ciò comprende quadrati e aquiloni non quadrati, per esempio quelli ottenuti da due triangoli con cateti $2,3$.
}
